\documentclass[11pt,letterpaper]{article}

\usepackage[utf8]{inputenc}
\usepackage[spanish]{babel}
\usepackage[margin=2.5cm]{geometry}
\usepackage{amsmath,amssymb,amsthm}
\usepackage{enumitem}
\usepackage{titlesec}
\usepackage{hyperref}
\usepackage{xcolor}

\definecolor{unicaucaverde}{RGB}{0,90,60}
\hypersetup{colorlinks=true,linkcolor=unicaucaverde,urlcolor=unicaucaverde}

\titleformat{\section}{\normalfont\Large\bfseries\color{unicaucaverde}}{\thesection}{1em}{}
\titleformat{\subsection}{\normalfont\large\bfseries}{\thesubsection}{1em}{}

\theoremstyle{definition}
\newtheorem{definicion}{Definición}[section]
\newtheorem{ejemplo}[definicion]{Ejemplo}
\newtheorem{propiedad}[definicion]{Propiedad}

\newcommand{\R}{\mathbb{R}}

\title{\textcolor{unicaucaverde}{\Large Notas de clase --- Semana 1}\\[4pt]
  \large Sistemas de ecuaciones lineales: introducción y el caso de dos incógnitas}
\author{Álgebra Lineal para Ingeniería --- Universidad del Cauca}
\date{2026}

\begin{document}
\maketitle
\tableofcontents

\section{Introducción a los sistemas de ecuaciones lineales}

\begin{definicion}[Ecuación lineal]
Una ecuación lineal en las variables $x_1,\dots,x_n$ es una expresión de la forma
\[
  a_1 x_1 + a_2 x_2 + \cdots + a_n x_n = b,
\]
donde $a_1,\dots,a_n$ (los \emph{coeficientes}) y $b$ (el \emph{término independiente}) son
números reales (o complejos) dados. Ninguna variable aparece elevada a una potencia distinta de
1, ni multiplicada por otra variable, ni dentro de una función no lineal.
\end{definicion}

\begin{definicion}[Sistema de $m$ ecuaciones lineales con $n$ incógnitas]
Un sistema de $m$ ecuaciones lineales con $n$ incógnitas $x_1,\dots,x_n$ es un conjunto de $m$
ecuaciones lineales que deben satisfacerse simultáneamente:
\[
\begin{aligned}
  a_{11}x_1 + a_{12}x_2 + \cdots + a_{1n}x_n &= b_1 \\
  a_{21}x_1 + a_{22}x_2 + \cdots + a_{2n}x_n &= b_2 \\
  &\;\vdots \\
  a_{m1}x_1 + a_{m2}x_2 + \cdots + a_{mn}x_n &= b_m
\end{aligned}
\]
Los números $a_{ij}$ son los \emph{coeficientes} (fila $i$, columna $j$), y $b_i$ los términos
independientes. Una \emph{solución} del sistema es una $n$-tupla $(x_1,\dots,x_n)$ que satisface
las $m$ ecuaciones a la vez. El \emph{conjunto solución} es el conjunto de todas las soluciones.
\end{definicion}

\begin{definicion}[Sistemas consistentes e inconsistentes]
Un sistema es \emph{consistente} si tiene al menos una solución, e \emph{inconsistente} si no
tiene ninguna. Un sistema consistente tiene exactamente una solución (\emph{determinado}) o
infinitas soluciones (\emph{indeterminado}) --- nunca un número finito mayor que uno; esto se
demuestra formalmente en la Semana 2, pero ya se observa en el caso $2\times 2$ de la Sección 2.
\end{definicion}

En ingeniería, cada ecuación del sistema suele representar una ley física o una restricción: una
ley de Kirchhoff de corrientes en un nodo, una ecuación de equilibrio de fuerzas en un punto de
una armadura, un balance de materia en una etapa de un proceso. El sistema completo describe el
estado del sistema físico bajo todas las restricciones simultáneamente.

\begin{ejemplo}
El sistema
\[
\begin{aligned}
  2x_1 - x_2 + 3x_3 &= 5 \\
  x_1 + x_2 - x_3 &= 0
\end{aligned}
\]
tiene $m=2$ ecuaciones y $n=3$ incógnitas. La tupla $(1,-1,-2)$ \emph{no} es solución: al
sustituir en la primera ecuación, $2(1)-(-1)+3(-2) = 2+1-6=-3 \neq 5$. En cambio, $(1,-1,0)$
tampoco lo es (verificarlo es un ejercicio); encontrar una solución real requiere el método
sistemático de la Semana 2 (Gauss-Jordan), porque aquí $n>m$ y hay más de una incógnita libre.
\end{ejemplo}

\section{Dos ecuaciones lineales con dos incógnitas}

\begin{definicion}[Sistema $2\times 2$]
Un sistema de dos ecuaciones lineales con dos incógnitas $x,y$ tiene la forma
\[
\begin{aligned}
  a_1 x + b_1 y &= c_1 \\
  a_2 x + b_2 y &= c_2
\end{aligned}
\]
Cada ecuación, si $(a_i,b_i)\neq(0,0)$, es la ecuación cartesiana de una recta en $\R^2$. Resolver
el sistema es hallar los puntos $(x,y)$ que están en \emph{ambas} rectas a la vez: la intersección.
\end{definicion}

\begin{propiedad}[Los tres casos]
Geométricamente, dos rectas en el plano solo pueden relacionarse de tres formas, que corresponden
exactamente a los tres tipos de sistema:
\begin{enumerate}[label=(\alph*)]
  \item \textbf{Rectas secantes} (se cortan en un punto) $\Rightarrow$ \textbf{solución única}.
  \item \textbf{Rectas paralelas distintas} (misma pendiente, no coinciden) $\Rightarrow$
        \textbf{sin solución} (sistema inconsistente).
  \item \textbf{Rectas coincidentes} (la misma recta escrita dos veces, salvo múltiplo escalar)
        $\Rightarrow$ \textbf{infinitas soluciones} (todos los puntos de la recta).
\end{enumerate}
\end{propiedad}

\begin{definicion}[Método de sustitución]
Se despeja una variable de una ecuación y se sustituye en la otra, reduciendo el sistema a una
sola ecuación con una incógnita.
\end{definicion}

\begin{definicion}[Método de eliminación]
Se multiplican las ecuaciones por constantes adecuadas y se suman (o restan) para cancelar una de
las variables, reduciendo también a una ecuación con una incógnita. Este método generaliza
directamente a la eliminación de Gauss-Jordan de la Semana 2.
\end{definicion}

\begin{ejemplo}[Solución única]
\[
\begin{aligned}
  x + y &= 5 \\
  x - y &= 1
\end{aligned}
\]
Sumando ambas ecuaciones: $2x = 6 \Rightarrow x=3$; sustituyendo, $y=2$. Solución única: $(3,2)$.
Las rectas $x+y=5$ y $x-y=1$ tienen pendientes distintas ($-1$ y $1$), así que se cortan en un
único punto.
\end{ejemplo}

\begin{ejemplo}[Sin solución]
\[
\begin{aligned}
  2x + y &= 4 \\
  4x + 2y &= 3
\end{aligned}
\]
La segunda ecuación es el doble de $2x+y$ igualado a $3$ en vez de $8$: mismas pendientes, mismo
"lado" pero distinto término independiente $\Rightarrow$ rectas paralelas distintas. El sistema
es inconsistente: no hay $(x,y)$ que satisfaga ambas.
\end{ejemplo}

\begin{ejemplo}[Infinitas soluciones]
\[
\begin{aligned}
  x - 2y &= 3 \\
  -2x + 4y &= -6
\end{aligned}
\]
La segunda ecuación es $-2$ veces la primera: es la misma recta. Todo punto de $x-2y=3$ es
solución; el conjunto solución es $\{(x,y) : x = 3+2y,\; y\in\R\}$, infinito.
\end{ejemplo}


\section{Soluciones de los ejercicios propuestos}

\begin{description}

\item[\textbf{Ejercicio 1.}]
Forma general de un sistema de 3 ecuaciones ($m=3$) con 4 incógnitas ($n=4$, $x_1, x_2, x_3, x_4$):
\[
\begin{aligned}
  a_{11}x_1 + a_{12}x_2 + a_{13}x_3 + a_{14}x_4 &= b_1 \\
  a_{21}x_1 + a_{22}x_2 + a_{23}x_3 + a_{24}x_4 &= b_2 \\
  a_{31}x_1 + a_{32}x_2 + a_{33}x_3 + a_{34}x_4 &= b_3
\end{aligned}
\]
donde $a_{ij} \in \R$ ($i=1,2,3;\; j=1,2,3,4$) son los coeficientes y $b_1, b_2, b_3 \in \R$ los términos independientes.

\item[\textbf{Ejercicio 2.}]
Para probar si $(x_1,x_2,x_3) = (2,-1,1)$ es solución del sistema:
\begin{enumerate}[label=(\roman*)]
  \item Ecuación 1: $x_1+x_2+x_3 = 2 + (-1) + 1 = 2$. Satisface la primera ecuación ($2=2$).
  \item Ecuación 2: $2x_1-x_2+x_3 = 2(2) - (-1) + 1 = 4 + 1 + 1 = 6$. Satisface la segunda ecuación ($6=6$).
\end{enumerate}
\[\boxed{\text{Sí es solución del sistema.}}\]

\item[\textbf{Ejercicio 3.}]
Simplificando las ecuaciones del circuito:
\begin{align*}
  10 = 4I_1 + 2(I_1-I_2) &\implies 6I_1 - 2I_2 = 10 \implies 3I_1 - I_2 = 5 \\
  0 = 2(I_2-I_1) + 3I_2 &\implies -2I_1 + 5I_2 = 0
\end{align*}
Forma estándar:
\[
\begin{aligned}
  3I_1 - I_2 &= 5 \\
  -2I_1 + 5I_2 &= 0
\end{aligned}
\]
Resolviendo por eliminación: multiplicamos la Ecuación 1 por 5 y sumamos con la Ecuación 2:
\[
(15I_1 - 5I_2) + (-2I_1 + 5I_2) = 25 + 0 \implies 13I_1 = 25 \implies I_1 = \frac{25}{13} \text{ A.}
\]
Sustituyendo $I_1$ en la Ecuación 2: $5I_2 = 2I_1 \implies 5I_2 = \frac{50}{13} \implies I_2 = \frac{10}{13} \text{ A.}$
\[\boxed{I_1 = \frac{25}{13}\text{ A},\quad I_2 = \frac{10}{13}\text{ A.}}\]

\item[\textbf{Ejercicio 4.}]
Sistema $\{3x-y=7,\; x+2y=0\}$.
De la segunda ecuación: $x = -2y$. Sustituyendo en la primera:
\[
3(-2y) - y = 7 \implies -7y = 7 \implies y = -1.
\]
Entonces $x = -2(-1) = 2$.
\[\boxed{\text{Solución única: }(2,-1).\quad \text{Sistema consistente determinado.}}\]

\item[\textbf{Ejercicio 5.}]
Sistema $\{x+3y=6,\; 2x+6y=10\}$.
Las pendientes son iguales: $m_1 = -1/3$, $m_2 = -2/6 = -1/3$.
Comparando razones de coeficientes: $\frac{1}{2} = \frac{3}{6} \neq \frac{6}{10}$.
Las rectas son paralelas distintas.
\[\boxed{\text{El sistema \textbf{no tiene solución} (inconsistente).}}\]

\item[\textbf{Ejercicio 6.}]
Dada la primera ecuación $3x+y=9$, una segunda ecuación que genera infinitas soluciones debe ser un múltiplo no nulo de la primera, por ejemplo multiplicando por $2$:
\[\boxed{6x + 2y = 18.}\]
El sistema $\{3x+y=9,\; 6x+2y=18\}$ representa rectas coincidentes y tiene infinitas soluciones.

\end{description}

\end{document}
